\documentclass{amsart}
% For dealing with references we use the comment environment
\usepackage{verbatim}
\newenvironment{reference}{\comment}{\endcomment}
%\newenvironment{reference}{}{}
\newenvironment{slogan}{\comment}{\endcomment}
\newenvironment{history}{\comment}{\endcomment}

% For commutative diagrams we use Xy-pic
\usepackage[all]{xy}

% We use 2cell for 2-commutative diagrams.
\xyoption{2cell}
\UseAllTwocells

% We use multicol for the list of chapters between chapters
\usepackage{multicol}

% This is generally recommended for better output
\usepackage{lmodern}
\usepackage[T1]{fontenc}

% For cross-file-references

% Package for hypertext links:
\usepackage{hyperref}

% For any local file, say "hello.tex" you want to link to please

% Theorem environments.
%
\theoremstyle{plain}
\newtheorem{theorem}[subsection]{Theorem}
\newtheorem{proposition}[subsection]{Proposition}
\newtheorem{lemma}[subsection]{Lemma}

\theoremstyle{definition}
\newtheorem{definition}[subsection]{Definition}
\newtheorem{example}[subsection]{Example}
\newtheorem{exercise}[subsection]{Exercise}
\newtheorem{situation}[subsection]{Situation}

\theoremstyle{remark}
\newtheorem{remark}[subsection]{Remark}
\newtheorem{remarks}[subsection]{Remarks}

\numberwithin{equation}{subsection}

% Macros
%
\def\lim{\mathop{\mathrm{lim}}\nolimits}
\def\colim{\mathop{\mathrm{colim}}\nolimits}
\def\Spec{\mathop{\mathrm{Spec}}}
\def\Hom{\mathop{\mathrm{Hom}}\nolimits}
\def\Ext{\mathop{\mathrm{Ext}}\nolimits}
\def\SheafHom{\mathop{\mathcal{H}\!\mathit{om}}\nolimits}
\def\SheafExt{\mathop{\mathcal{E}\!\mathit{xt}}\nolimits}
\def\Sch{\mathit{Sch}}
\def\Mor{\mathop{\mathrm{Mor}}\nolimits}
\def\Ob{\mathop{\mathrm{Ob}}\nolimits}
\def\Sh{\mathop{\mathit{Sh}}\nolimits}
\def\NL{\mathop{N\!L}\nolimits}
\def\CH{\mathop{\mathrm{CH}}\nolimits}
\def\proetale{{pro\text{-}\acute{e}tale}}
\def\etale{{\acute{e}tale}}
\def\QCoh{\mathit{QCoh}}
\def\Ker{\mathop{\mathrm{Ker}}}
\def\Im{\mathop{\mathrm{Im}}}
\def\Coker{\mathop{\mathrm{Coker}}}
\def\Coim{\mathop{\mathrm{Coim}}}

% Boxtimes
%
\DeclareMathSymbol{\boxtimes}{\mathbin}{AMSa}{"02}

%
% Macros for moduli stacks/spaces
%
\def\QCohstack{\mathcal{QC}\!\mathit{oh}}
\def\Cohstack{\mathcal{C}\!\mathit{oh}}
\def\Spacesstack{\mathcal{S}\!\mathit{paces}}
\def\Quotfunctor{\mathrm{Quot}}
\def\Hilbfunctor{\mathrm{Hilb}}
\def\Curvesstack{\mathcal{C}\!\mathit{urves}}
\def\Polarizedstack{\mathcal{P}\!\mathit{olarized}}
\def\Complexesstack{\mathcal{C}\!\mathit{omplexes}}
% \Pic is the operator that assigns to X its picard group, usage \Pic(X)
% \Picardstack_{X/B} denotes the Picard stack of X over B
% \Picardfunctor_{X/B} denotes the Picard functor of X over B
\def\Pic{\mathop{\mathrm{Pic}}\nolimits}
\def\Picardstack{\mathcal{P}\!\mathit{ic}}
\def\Picardfunctor{\mathrm{Pic}}
\def\Deformationcategory{\mathcal{D}\!\mathit{ef}}

\setlength{\emergencystretch}{3em}
\begin{document}
\title[Stacks Project, Tag 0F1H]{289 --- Product base change for open immersions with invertible torsion coefficients}
\maketitle
\noindent Copyright (C) 2005--2025 Johan de Jong. Stacks Project authors. Source: \href{https://stacks.math.columbia.edu/tag/0F1H}{Tag 0F1H}.

\noindent Editorial difficulty note (not author text). Difficulty H3: The proof must handle an arbitrary nonsmooth product factor, so smooth base change alone is insufficient. A minimal-dimensional counterexample is forced onto closed points by strict-henselian dimension reduction; projective compactification and the one-proper-factor Kunneth formula then annihilate its cone..

\noindent Editorial provenance note (not author text). History: excerpt from revision \texttt{a0444\allowbreak 6e57e\allowbreak c1fbc\allowbreak 252a8\allowbreak 71afc\allowbreak ec775\allowbreak 2fb28\allowbreak 07b14}, etale-cohomology.tex, lines 20402--20594. Reference presentation was adapted; original-author context and core support blocks were retained or added. The mathematical wording is unchanged. Licensed under GNU FDL 1.2 or later; no invariant sections or cover texts. See ../licenses/COPYING. Context blocks reproduce author definitions and supporting results; they are not additional counted problems.

\begin{lemma}
\label{lemma-kunneth-localize-on-X}
Let $K$ be a field. Let $j : U \to X$ be an open immersion of
schemes of finite type over $K$. Let $Y$ be a scheme of finite type
over $K$. Consider the diagram
$$
\xymatrix{
Y \times_{\Spec(K)} X \ar[d]_q  &
Y \times_{\Spec(K)} U \ar[l]^h \ar[d]^p \\
X & U \ar[l]_j
}
$$
Then the base change map $q^{-1}Rj_*\mathcal{F} \to Rh_*p^{-1}\mathcal{F}$
is an isomorphism for $\mathcal{F}$ an abelian sheaf on $U_\etale$
whose stalks are torsion of orders invertible in $K$.
\end{lemma}

\begin{proof}
Write $\mathcal{F} = \colim \mathcal{F}[n]$ where the colimit
is over the multiplicative system of integers invertible in $K$.
Since cohomology commutes with filtered colimits in our situation
(for a precise reference see Lemma \href{https://stacks.math.columbia.edu/tag/0EZT}{lemma base change Rf star colim (Tag 0EZT)}),
it suffices to prove the lemma for $\mathcal{F}[n]$. Thus we may
assume $\mathcal{F}$ is a sheaf of $\mathbf{Z}/n\mathbf{Z}$-modules
for some $n$ invertible in $K$ (we will use this at the very end of
the proof). In the proof we use the short hand $X \times_K Y$ for the fibre
product over $\Spec(K)$.
We will prove the lemma by induction on $\dim(X) + \dim(Y)$.
The lemma is trivial if $\dim(X) \leq 0$, since in this case
$U$ is an open and closed subscheme of $X$.
Choose a point $z \in X \times_K Y$. We will show
the stalk at $\overline{z}$ is an isomorphism.

\medskip\noindent
Suppose that $z \mapsto x \in X$ and assume $\text{trdeg}_K(\kappa(x)) > 0$.
Set $X' = \Spec(\mathcal{O}_{X, x}^{sh})$
and denote $U' \subset X'$ the inverse image of $U$.
Consider the base change
$$
\xymatrix{
Y \times_K X' \ar[d]_{q'}  &
Y \times_K U' \ar[l]^{h'} \ar[d]^{p'} \\
X' & U' \ar[l]_{j'}
}
$$
of our diagram by $X' \to X$. Observe that $X' \to X$ is a filtered
colimit of \'etale morphisms. By smooth base change in the form of
Lemma \href{https://stacks.math.columbia.edu/tag/0F09}{lemma smooth base change general (Tag 0F09)} the pullback of
$q^{-1}Rj_*\mathcal{F} \to Rh_*p^{-1}\mathcal{F}$ to $X'$ to
$Y \times_K X'$ is the map
$(q')^{-1}Rj'_*\mathcal{F}' \to Rj'_*(p')^{-1}\mathcal{F}'$ where
$\mathcal{F}'$ is the pullback of $\mathcal{F}$ to $U'$.
(In this step it would suffice to use \'etale base change which is
an essentially trivial result.) So it suffices to show
that $(q')^{-1}Rj'_*\mathcal{F}' \to Rj'_*(p')^{-1}\mathcal{F}'$
is an isomorphism in order to prove that our original
map is an isomorphism on stalks at $\overline{z}$.
By Lemma \href{https://stacks.math.columbia.edu/tag/0F0U}{lemma strictly henselian (Tag 0F0U)}
there is a separably closed field $L/K$ such that
$X' = \lim X_i$ with $X_i$ affine of finite type over $L$
and $\dim(X_i) < \dim(X)$. For $i$ large enough there
exists an open $U_i \subset X_i$ restricting to $U'$ in $X'$.
We may apply the induction hypothesis to the diagram
$$
\vcenter{
\xymatrix{
Y \times_K X_i \ar[d]_{q_i}  &
Y \times_K U_i \ar[l]^{h_i} \ar[d]^{p_i} \\
X_i & U_i \ar[l]_{j_i}
}
}
\quad\text{equal to}\quad
\vcenter{
\xymatrix{
Y_L \times_L X_i \ar[d]_{q_i}  &
Y_L \times_L U_i \ar[l]^{h_i} \ar[d]^{p_i} \\
X_i & U_i \ar[l]_{j_i}
}
}
$$
over the field $L$ and the pullback of $\mathcal{F}$ to these diagrams.
By Lemma \href{https://stacks.math.columbia.edu/tag/0EZT}{lemma base change Rf star colim (Tag 0EZT)} we conclude that the map
$(q')^{-1}Rj'_*\mathcal{F}' \to Rj'_*(p')^{-1}\mathcal{F}$ is an isomorphism.

\medskip\noindent
Suppose that $z \mapsto y \in Y$ and assume $\text{trdeg}_K(\kappa(y)) > 0$.
Let $Y' = \Spec(\mathcal{O}_{X, x}^{sh})$.
By Lemma \href{https://stacks.math.columbia.edu/tag/0F0U}{lemma strictly henselian (Tag 0F0U)} there is a separably closed field
$L/K$ such that $Y' = \lim Y_i$ with $Y_i$ affine of finite type over $L$
and $\dim(Y_i) < \dim(Y)$. In particular $Y'$ is a scheme over $L$.
Denote with a subscript $L$ the base change from schemes over $K$
to schemes over $L$. Consider the commutative diagrams
$$
\vcenter{
\xymatrix{
Y' \times_K X \ar[d]_f  &
Y' \times_K U \ar[l]^{h'} \ar[d]^{f'} \\
Y \times_K X \ar[d]_q  &
Y \times_K U \ar[l]^h \ar[d]^p \\
X & U \ar[l]_j
}
}
\quad\text{and}\quad
\vcenter{
\xymatrix{
Y' \times_L X_L \ar[d]_{q'}  &
Y' \times_L U_L \ar[l]^{h'} \ar[d]^{p'} \\
X_L \ar[d] &
U_L \ar[l]^{j_L} \ar[d] \\
X & U \ar[l]_j
}
}
$$
and observe the top and bottom rows are the same on the left and the right.
By smooth base change we see that
$f^{-1}Rh_*p^{-1}\mathcal{F} = Rh'_*(f')^{-1}p^{-1}\mathcal{F}$
(similarly to the previous paragraph).
By smooth base change for $\Spec(L) \to \Spec(K)$
(Lemma \href{https://stacks.math.columbia.edu/tag/0F1C}{lemma base change field extension (Tag 0F1C)})
we see that $Rj_{L, *}\mathcal{F}_L$ is the pullback of
$Rj_*\mathcal{F}$ to $X_L$. Combining these two observations,
we conclude that it suffices to prove the base change map
for the upper square in the diagram on the right
is an isomorphism in order to prove that our original
map is an isomorphism on stalks at $\overline{z}$\footnote{Here we
use that a ``vertical composition'' of base change maps is a base change
map as explained in Cohomology on Sites, Remark
\href{https://stacks.math.columbia.edu/tag/0E46}{remark compose base change (Tag 0E46)}.}.
Then using that $Y' = \lim Y_i$ and argueing exactly
as in the previous paragraph we see that the induction
hypothesis forces our map over $Y' \times_K X$ to be
an isomorphism.

\medskip\noindent
Thus any counter example with $\dim(X) + \dim(Y)$ minimal
would only have nonisomorphisms
$q^{-1}Rj_*\mathcal{F} \to Rh_*p^{-1}\mathcal{F}$
on stalks at closed points of $X \times_K Y$ (because a point
$z$ of $X \times_K Y$ is a closed point if and only if
both the image of $z$ in $X$ and in $Y$ are closed).
Since it is enough to prove the isomorphism locally,
we may assume $X$ and $Y$ are affine. However, then
we can choose an open dense immersion $Y \to Y'$ with $Y'$
projective. (Choose a closed immersion $Y \to \mathbf{A}^n_K$
and let $Y'$ be the scheme theoretic closure of $Y$ in
$\mathbf{P}^n_K$.) Then $\dim(Y') = \dim(Y)$ and hence
we get a ``minimal'' counter example with $Y$ projective over $K$.
In the next paragraph we show that this can't happen.

\medskip\noindent
Consider a diagram as in the statement of the lemma
such that $q^{-1}Rj_*\mathcal{F} \to Rh_*p^{-1}\mathcal{F}$
is an isomorphism at all non-closed points of $X \times_K Y$
and such that $Y$ is projective.
The restriction of the map to $(X \times_K Y)_{K^{sep}}$
is the corresponding map for the diagram of the lemma
base changed to $K^{sep}$. Thus we may and do assume $K$
is separably algebraically closed. Choose a distinguished triangle
$$
q^{-1}Rj_*\mathcal{F} \to Rh_*p^{-1}\mathcal{F} \to Q \to
(q^{-1}Rj_*\mathcal{F})[1]
$$
in $D((X \times_K Y)_\etale)$. Since $Q$ is supported in
closed points we see that it suffices to prove
$H^i(X \times_K Y, Q) = 0$ for all $i$, see
Lemma \href{https://stacks.math.columbia.edu/tag/0F1G}{lemma vanishing closed points (Tag 0F1G)}.
Thus it suffices to prove that
$q^{-1}Rj_*\mathcal{F} \to Rh_*p^{-1}\mathcal{F}$
induces an isomorphism on cohomology. Recall that $\mathcal{F}$
is annihilated by $n$ invertible in $K$.
By the K\"unneth formula of Lemma \href{https://stacks.math.columbia.edu/tag/0F14}{lemma kunneth one proper (Tag 0F14)}
we have
\begin{align*}
R\Gamma(X \times_K Y, q^{-1}Rj_*\mathcal{F})
&=
R\Gamma(X, Rj_*\mathcal{F})
\otimes_{\mathbf{Z}/n\mathbf{Z}}^\mathbf{L}
R\Gamma(Y, \mathbf{Z}/n\mathbf{Z}) \\
& =
R\Gamma(U, \mathcal{F})
\otimes_{\mathbf{Z}/n\mathbf{Z}}^\mathbf{L}
R\Gamma(Y, \mathbf{Z}/n\mathbf{Z})
\end{align*}
and
$$
R\Gamma(X \times_K Y, Rh_*p^{-1}\mathcal{F}) =
R\Gamma(U \times_K Y, p^{-1}\mathcal{F}) =
R\Gamma(U, \mathcal{F})
\otimes_{\mathbf{Z}/n\mathbf{Z}}^\mathbf{L}
R\Gamma(Y, \mathbf{Z}/n\mathbf{Z})
$$
This finishes the proof.
\end{proof}
\end{document}
